\UseRawInputEncoding
\documentclass[aspectratio=169, 10pt, xcolor=table]{beamer}

% --- Configuracion del Tema ---
\usetheme{Madrid}
\usecolortheme{dolphin}

% --- Paquetes necesarios ---
\usepackage[utf8]{inputenc}
\usepackage[T1]{fontenc}
\usepackage[spanish,es-noquoting]{babel}
\usepackage{amsmath, amssymb}
\usepackage{mathtools}
\usepackage{booktabs}
\usepackage{graphicx}
\usepackage{listings}
\usepackage{tikz}
\usetikzlibrary{arrows.meta, positioning, shapes.geometric, babel}
\usepackage{etoolbox}
\usepackage{upquote}

% --- Estilo para codigo SQL ---
\definecolor{codegreen}{rgb}{0,0.6,0}
\definecolor{codegray}{rgb}{0.5,0.5,0.5}
\definecolor{codepurple}{rgb}{0.58,0,0.82}
\definecolor{backcolour}{rgb}{0.95,0.95,0.92}
\lstset{
    language=SQL,
    backgroundcolor=\color{backcolour},
    commentstyle=\color{codegreen},
    keywordstyle=\color{blue},
    numberstyle=\tiny\color{codegray},
    stringstyle=\color{codepurple},
    basicstyle=\ttfamily\scriptsize,
    breaklines=true,
    frame=single,
    showstringspaces=false,
    literate={á}{{\'a}}1 {é}{{\'e}}1 {í}{{\'i}}1 {ó}{{\'o}}1 {ú}{{\'u}}1 {ñ}{{\~n}}1 {ü}{{\"u}}1
}

% --- Pie de pagina personalizado ---
\makeatletter
\setbeamertemplate{footline}{
  \leavevmode%
  \hbox{%
  \begin{beamercolorbox}[wd=.333333\paperwidth,ht=2.25ex,dp=1ex,center]{author in head/foot}%
    \usebeamerfont{author in head/foot}\insertshortauthor
  \end{beamercolorbox}%
  \begin{beamercolorbox}[wd=.333333\paperwidth,ht=2.25ex,dp=1ex,center]{title in head/foot}%
    \usebeamerfont{title in head/foot}Sesión 4
  \end{beamercolorbox}%
  \begin{beamercolorbox}[wd=.333333\paperwidth,ht=2.25ex,dp=1ex,right]{date in head/foot}%
    \usebeamerfont{date in head/foot}BBDD \hfill \insertframenumber/\inserttotalframenumber \hspace*{0.3cm}
  \end{beamercolorbox}}%
  \vskip0pt%
}
\makeatother

\title[SESIÓN 4 BBDD]{\textbf{Sesión 4}}
\subtitle{Arquitectura, Optimización y Modelamiento Avanzado}
\author[Prof. C. Sánchez]{Carlos César Sánchez Coronel}
\institute{\textbf{BASES DE DATOS}}
\date{}

\setbeamertemplate{navigation symbols}{}

\AtBeginSection[]{
  \begin{frame}
    \frametitle{Contenido de la sección}
    \tableofcontents[currentsection]
  \end{frame}
}

\begin{document}

\begin{frame}
    \titlepage
\end{frame}

\begin{frame}{Contenido general}
    \tableofcontents
\end{frame}

% ============================================================================
% SECCIÓN 1: ARQUITECTURAS DE BASES DE DATOS
% ============================================================================
\section{Arquitecturas de Bases de Datos}

\begin{frame}{Modelos Clásicos}
    \begin{columns}[t]
        \column{0.5\textwidth}
        \textbf{Centralizado}
        \begin{itemize}
            \item Único servidor.
            \item Simple pero punto único de fallo.
            \item Escalabilidad vertical.
        \end{itemize}
        \column{0.5\textwidth}
        \textbf{Cliente-Servidor}
        \begin{itemize}
            \item Separación de responsabilidades.
            \item Estándar en SGBD (PostgreSQL, MySQL, Oracle).
            \item Escalabilidad limitada.
        \end{itemize}
    \end{columns}
\end{frame}

\begin{frame}{Bases de Datos Distribuidas}
    \textbf{Definición:} Datos almacenados en múltiples nodos interconectados.
    \begin{block}{Técnicas principales}
        \begin{itemize}
            \item \textbf{Fragmentación}: Horizontal (filas), Vertical (columnas), Híbrida.
            \item \textbf{Replicación}: Copias en varios servidores (síncrona/asíncrona).
            \item \textbf{Transparencia}: El usuario no percibe la distribución.
        \end{itemize}
    \end{block}
    \begin{exampleblock}{Ejemplo fragmentación horizontal}
        Clientes de Europa en un nodo, América en otro.
    \end{exampleblock}
\end{frame}

\begin{frame}{Alta Disponibilidad y Sharding}
    \begin{columns}[t]
        \column{0.5\textwidth}
        \textbf{Clustering}
        \begin{itemize}
            \item Active-Passive (standby)
            \item Active-Active (todos activos)
            \item Objetivo: alta disponibilidad
        \end{itemize}
        \column{0.5\textwidth}
        \textbf{Sharding}
        \begin{itemize}
            \item Particionamiento horizontal masivo.
            \item Cada \textit{shard} es una BD independiente.
            \item Clave de partición crucial.
            \item Ej: redes sociales.
        \end{itemize}
    \end{columns}
    \vspace{0.3cm}
    \centering
    \begin{tikzpicture}[scale=0.6, every node/.style={scale=0.8}]
        \draw[fill=blue!20] (0,0) rectangle (2,1) node[midway] {App};
        \draw[->] (2,0.5) -- (3,0.5);
        \draw[fill=green!20] (3,1) rectangle (5,2) node[midway, align=center] {Shard 1\\ ids 1-1000};
        \draw[fill=green!20] (3,-0.5) rectangle (5,0.5) node[midway, align=center] {Shard 2\\ ids 1001-2000};
        \draw[fill=green!20] (3,-2) rectangle (5,-1) node[midway, align=center] {Shard 3\\ ids 2001-3000};
    \end{tikzpicture}
\end{frame}

\begin{frame}{Bases de Datos en la Nube (DBaaS)}
    \textbf{Plataforma como servicio:}
    \begin{itemize}
        \item Amazon RDS, Google Cloud SQL, Azure SQL, MongoDB Atlas.
    \end{itemize}
    \begin{block}{Ventajas}
        \begin{itemize}
            \item Escalabilidad automática.
            \item Backups y parches gestionados.
            \item Alta disponibilidad integrada.
            \item Pago por uso.
        \end{itemize}
    \end{block}
\end{frame}

% ============================================================================
% SECCIÓN 2: OPTIMIZACIÓN DE CONSULTAS E ÍNDICES
% ============================================================================
\section{Optimización de Consultas e Índices}

\begin{frame}[fragile]{El Optimizador y los Planes de Ejecución}
    \begin{itemize}
        \item \textbf{Optimizador}: Decide el plan más eficiente (basado en costos).
    \end{itemize}
    \begin{block}{Operadores comunes}
        \begin{description}
            \item[Seq Scan:] Recorre toda la tabla (O(n)).
            \item[Index Scan:] Usa un índice (O(log n)).
            \item[Index Only Scan:] Solo el índice, sin tabla.
            \item[Hash Join:] Tabla hash en memoria.
        \end{description}
    \end{block}
\begin{verbatim}
EXPLAIN (ANALYZE, BUFFERS) 
SELECT * FROM clientes WHERE ciudad = 'Lima';
\end{verbatim}
\end{frame}

\begin{frame}{Estadísticas y Tipos de Índices Avanzados}
    \begin{itemize}
        \item \textbf{Estadísticas}: Número de filas, histogramas. Se actualizan con \texttt{ANALYZE}.
    \end{itemize}
    \begin{columns}[t]
        \column{0.33\textwidth}
        \textbf{Índice Cubriente}
        \begin{itemize}
            \item Incluye columnas extra.
            \item Evita acceso a tabla.
        \end{itemize}
        \column{0.33\textwidth}
        \textbf{Índice Funcional}
        \begin{itemize}
            \item Indexa una función.
            \item Ej: \texttt{LOWER(email)}.
        \end{itemize}
        \column{0.34\textwidth}
        \textbf{Índice Parcial}
        \begin{itemize}
            \item Solo un subconjunto.
            \item Ej: \texttt{WHERE activo = true}.
        \end{itemize}
    \end{columns}
\end{frame}

\begin{frame}[fragile]{Ejercicio Resuelto: Optimización de Consulta Lenta}
    \textbf{Problema:}
\begin{verbatim}
SELECT * FROM ventas 
WHERE fecha BETWEEN '2025-01-01' AND '2025-01-31' 
ORDER BY total DESC;
\end{verbatim}
    Tabla de 10M filas, sin índices $\rightarrow$ Seq Scan + Sort (lento).
    
    \textbf{Solución:}
\begin{verbatim}
CREATE INDEX idx_ventas_fecha_total 
ON ventas (fecha, total DESC);
\end{verbatim}
    \begin{itemize}
        \item Index Scan por fecha.
        \item Datos ya ordenados (evita Sort).
    \end{itemize}
\end{frame}

% ============================================================================
% SECCIÓN 3: TRANSACCIONES Y CONCURRENCIA
% ============================================================================
\section{Transacciones y Concurrencia}

\begin{frame}{Propiedades ACID}
    \begin{description}
        \item[Atomicidad] Todo o nada (logs undo/redo).
        \item[Consistencia] Integridad de datos (restricciones).
        \item[Aislamiento] Transacciones no interfieren (locks / MVCC).
        \item[Durabilidad] Persiste tras fallo (WAL).
    \end{description}
\end{frame}

\begin{frame}{Niveles de Aislamiento (Estándar SQL)}
    \begin{table}
        \centering
        \begin{tabular}{l|c|c|c}
            \toprule
            Nivel & Dirty Read & Non-repeatable & Phantom \\
            \midrule
            Read Uncommitted & Posible & Posible & Posible \\
            Read Committed & No & Posible & Posible \\
            Repeatable Read & No & No & Posible \\
            Serializable & No & No & No \\
            \bottomrule
        \end{tabular}
    \end{table}
    \begin{itemize}
        \item \textbf{PostgreSQL}: Default \texttt{READ COMMITTED}.
    \end{itemize}
\end{frame}

\begin{frame}{Mecanismos de Concurrencia}
    \begin{block}{Bloqueos (Locks)}
        \begin{itemize}
            \item Compartidos (lectura), Exclusivos (escritura).
            \item Pueden causar \textbf{deadlocks}.
        \end{itemize}
    \end{block}
    \begin{block}{MVCC (Multi-Version Concurrency Control)}
        \begin{itemize}
            \item Cada transacción ve una \textit{instantánea}.
            \item Lecturas NO bloquean escrituras.
            \item Usado por PostgreSQL, Oracle.
        \end{itemize}
    \end{block}
\end{frame}

\begin{frame}[fragile]{Ejemplo de Concurrencia en PostgreSQL}
    \textbf{Dos sesiones actualizando la misma fila:}
\begin{verbatim}
-- Sesion 1
BEGIN;
UPDATE cuenta SET saldo = saldo - 100 WHERE id = 1;

-- Sesion 2
BEGIN;
UPDATE cuenta SET saldo = saldo + 100 WHERE id = 1; -- Bloqueada
\end{verbatim}
    \begin{itemize}
        \item En \texttt{READ COMMITTED}: T2 espera a T1.
        \item En \texttt{REPEATABLE READ}: Si T1 hace commit, T2 falla con error de serialización.
    \end{itemize}
\end{frame}

% ============================================================================
% SECCIÓN 4: PROGRAMACIÓN EN BASE DE DATOS
% ============================================================================
\section{Programación en Base de Datos}

% CORREGIDO: Se añade [fragile] para permitir verbatim
\begin{frame}[fragile]{Triggers y Restricciones Avanzadas}
    \textbf{Trigger:} Procedimiento automático ante un evento (INSERT, UPDATE, DELETE).
    \begin{itemize}
        \item \texttt{BEFORE}, \texttt{AFTER}, \texttt{INSTEAD OF}.
        \item Nivel: \texttt{FOR EACH ROW} o \texttt{FOR EACH STATEMENT}.
    \end{itemize}
    \begin{exampleblock}{Ejemplo: Trigger de auditoría de salario}
\begin{verbatim}
CREATE TRIGGER trig_audit_salario
AFTER UPDATE OF salario ON empleados
FOR EACH ROW EXECUTE FUNCTION audit_salario_func();
\end{verbatim}
    \end{exampleblock}
\end{frame}

\begin{frame}{Procedimientos, Funciones y Vistas Materializadas}
    \begin{columns}[t]
        \column{0.5\textwidth}
        \textbf{Funciones / Procedimientos}
        \begin{itemize}
            \item Encapsulan lógica.
            \item Reducen tráfico de red.
            \item Escritos en PL/pgSQL, T-SQL.
        \end{itemize}
        \column{0.5\textwidth}
        \textbf{Vistas Materializadas}
        \begin{itemize}
            \item Almacenan resultado físico.
            \item Útiles para reportes.
            \item Se refrescan (\texttt{REFRESH MATERIALIZED VIEW}).
        \end{itemize}
    \end{columns}
\end{frame}

% ============================================================================
% SECCIÓN 5: DATA WAREHOUSING Y MODELAMIENTO DIMENSIONAL
% ============================================================================
\section{Data Warehousing y Modelamiento Dimensional}

\begin{frame}{OLTP vs OLAP}
    \begin{table}
        \centering
        \begin{tabular}{l|l|l}
            \toprule
            Característica & OLTP & OLAP \\
            \midrule
            Propósito & Transacciones diarias & Análisis histórico \\
            Modelado & Normalizado (3FN) & Desnormalizado (Estrella) \\
            Consultas & Cortas, frecuentes & Complejas, agregaciones \\
            Optimizado para & Escritura & Lectura \\
            \bottomrule
        \end{tabular}
    \end{table}
\end{frame}

\begin{frame}{Esquema Estrella (Star Schema)}
    \begin{columns}[t]
        \column{0.3\textwidth}
        \textbf{Componentes:}
        \begin{itemize}
            \item Tabla de Hechos (medidas, FKs).
            \item Tablas de Dimensiones (descripciones).
        \end{itemize}
        \column{0.7\textwidth}
        \centering
        \begin{tikzpicture}[node distance=1.5cm, auto, every node/.style={draw, minimum width=2cm, minimum height=0.8cm, align=center}]
            \node (hechos) {hechos\_ventas};
            \node[above=of hechos] (tiempo) {dim\_tiempo};
            \node[left=of hechos] (cliente) {dim\_cliente};
            \node[right=of hechos] (sucursal) {dim\_sucursal};
            \node[below=of hechos] (producto) {dim\_producto};
            \draw[->] (tiempo) -- (hechos);
            \draw[->] (cliente) -- (hechos);
            \draw[->] (sucursal) -- (hechos);
            \draw[->] (producto) -- (hechos);
        \end{tikzpicture}
    \end{columns}
\end{frame}

\begin{frame}{Dimensiones Lentamente Cambiantes (SCD)}
    \textbf{Manejo de cambios en atributos de dimensiones:}
    \begin{description}
        \item[Tipo 1:] Sobrescribir. Sin histórico. (Ej: corregir error)
        \item[Tipo 2:] Crear nueva fila con fechas de vigencia. Histórico completo. (Ej: cambio de precio)
        \item[Tipo 3:] Agregar columnas para valor anterior/actual. Histórico limitado.
    \end{description}
\end{frame}

% ============================================================================
% SECCIÓN 6: INTRODUCCIÓN A NOSQL (SOLO TEORÍA)
% ============================================================================
\section{Introducción a NoSQL}

\begin{frame}{Clasificación de Bases de Datos NoSQL}
    \begin{table}
        \centering
        \begin{tabular}{l|l}
            \toprule
            Tipo & Ejemplos \\
            \midrule
            Clave-Valor & Redis, DynamoDB \\
            Documentales & MongoDB, Couchbase \\
            Columnares (Wide-Column) & Cassandra, HBase \\
            Grafos & Neo4j, Amazon Neptune \\
            \bottomrule
        \end{tabular}
    \end{table}
\end{frame}

\begin{frame}{Teorema CAP}
    \begin{block}{En sistemas distribuidos, solo dos de tres:}
        \begin{itemize}
            \item \textbf{C}onsistencia (todos ven lo mismo).
            \item \textbf{D}isponibilidad (siempre responde).
            \item \textbf{T}olerancia a Particiones (funciona pese a fallos de red).
        \end{itemize}
    \end{block}
    \begin{columns}[t]
        \column{0.5\textwidth}
        \textbf{CP (Consistencia + Partición)}
        \begin{itemize}
            \item HBase, MongoDB (típico)
        \end{itemize}
        \column{0.5\textwidth}
        \textbf{AP (Disponibilidad + Partición)}
        \begin{itemize}
            \item Cassandra, DynamoDB
        \end{itemize}
    \end{columns}
\end{frame}

\begin{frame}{Modelamiento en MongoDB (Conceptos)}
    \begin{itemize}
        \item \textbf{Documentos anidados}: datos relacionados se guardan juntos (desnormalización).
        \item \textbf{Referencias}: se guarda el ID de otro documento (similar a FK, sin integridad automática).
        \item Ejemplo: un pedido contiene un array de productos.
    \end{itemize}
\end{frame}

\begin{frame}{Modelamiento en Cassandra (Conceptos)}
    \begin{itemize}
        \item Modelado orientado a consultas.
        \item Clave primaria compuesta:
        \begin{itemize}
            \item \textbf{Clave de partición}: determina el nodo.
            \item \textbf{Columnas de clustering}: orden dentro de la partición.
        \end{itemize}
        \item Las consultas deben incluir la clave de partición para ser eficientes.
    \end{itemize}
\end{frame}

\begin{frame}{¿Cuándo usar SQL y cuándo NoSQL?}
    \begin{block}{Elegir SQL si...}
        \begin{itemize}
            \item Transacciones ACID estrictas.
            \item Consultas complejas con joins.
            \item Esquema estable e integridad referencial.
        \end{itemize}
    \end{block}
    \begin{block}{Elegir NoSQL si...}
        \begin{itemize}
            \item Escalabilidad horizontal masiva.
            \item Esquema flexible y cambiante.
            \item Se tolera consistencia eventual.
            \item Alto rendimiento para patrones simples.
        \end{itemize}
    \end{block}
\end{frame}

% ============================================================================
% SECCIÓN 7: CASOS PRÁCTICOS INTEGRADORES
% ============================================================================
\section{Casos Prácticos Integradores}

\begin{frame}{Ejercicio 1: Banca en Línea}
    \textbf{Requisitos:} Cuentas, clientes, transacciones, sucursales.
    \begin{itemize}
        \item \textbf{Concurrencia}: Usar transacciones con \texttt{READ COMMITTED} y \texttt{SELECT FOR UPDATE}.
        \item \textbf{Particionamiento}: Tabla \texttt{Transaccion} particionada por fecha.
        \item \textbf{Índices}: Índice en \texttt{Transaccion(id\_cuenta, fecha)}.
    \end{itemize}
\end{frame}

\begin{frame}[fragile]{Ejercicio 2: Data Warehouse de Ventas (SCD Tipo 2)}
    \begin{block}{Dimensión Producto con histórico de precios}
\begin{verbatim}
CREATE TABLE dim_producto (
    id_producto INT,
    nombre TEXT,
    precio DECIMAL,
    fecha_inicio DATE,
    fecha_fin DATE,
    activo BOOLEAN
);
\end{verbatim}
    \end{block}
    \begin{itemize}
        \item Cada cambio de precio inserta una nueva fila.
        \item La tabla de hechos apunta al \texttt{id\_producto} vigente.
    \end{itemize}
\end{frame}

\begin{frame}[fragile]{Ejercicio 3: Trigger de Histórico}
\begin{verbatim}
CREATE TABLE historial_empleado (...);
CREATE OR REPLACE FUNCTION historial_empleado_func() 
RETURNS TRIGGER AS $$
BEGIN
    INSERT INTO historial_empleado (id_empleado, salario_anterior, ...)
    VALUES (OLD.id, OLD.salario, NEW.salario, ...);
    RETURN NEW;
END;
$$ LANGUAGE plpgsql;
CREATE TRIGGER trig_historial_empleado
AFTER UPDATE OF salario, departamento_id ON empleados
FOR EACH ROW EXECUTE FUNCTION historial_empleado_func();
\end{verbatim}
\end{frame}

\begin{frame}{Ejercicio 4: Comparación SQL vs NoSQL para un Blog}
    \begin{columns}[t]
        \column{0.5\textwidth}
        \textbf{SQL}
        \begin{itemize}
            \item Tablas: Post, Comentario, Like (con FKs).
            \item Joins para obtener datos completos.
            \item Integridad referencial.
        \end{itemize}
        \column{0.5\textwidth}
        \textbf{MongoDB}
        \begin{itemize}
            \item Colección \texttt{posts} con array de comentarios embebidos.
            \item Una sola consulta.
            \item Límite de 16MB si hay muchos comentarios.
        \end{itemize}
    \end{columns}
\end{frame}

% ============================================================================
% SECCIÓN 8: RESUMEN
% ============================================================================
\section{Resumen}

\begin{frame}{Conceptos Clave de la Sesión 4}
    \begin{itemize}
        \item \textbf{Arquitecturas}: Centralizada, Distribuida, Sharding, DBaaS.
        \item \textbf{Optimización}: Planes de ejecución, índices (cubriente, funcional, parcial), estadísticas.
        \item \textbf{Transacciones}: ACID, niveles de aislamiento, MVCC, deadlocks.
        \item \textbf{Programación BD}: Triggers, procedimientos, vistas materializadas.
        \item \textbf{Data Warehousing}: OLTP/OLAP, Esquema Estrella, SCD, ETL.
        \item \textbf{NoSQL}: Tipos, Teorema CAP, Modelamiento conceptual.
        \item \textbf{Casos prácticos}: Banca, Ventas, Auditoría, Comparación SQL/NoSQL.
    \end{itemize}
\end{frame}

\begin{frame}
    \centering

    
    \vspace{1cm}
    \large ¡Gracias!
\end{frame}

\end{document}